\chapter{\texttt{pipeline.atmosphere\_earth} --- atmospheric surface to detector}
\label{ch:pipeline-atmosphere-earth}

\texttt{pipeline.atmosphere\_earth} continues a coherent atmospheric
surface state (\cref{ch:pipeline-atmosphere}) through the Earth to a
detector. It also hosts the detector-composition helpers that
\texttt{medium.atmosphere} deliberately does not carry itself
(\cref{ch:atmosphere}): the atmosphere medium only propagates production
$\to$ surface, while combining per-flavour detector probabilities with
production-flux tables, integrating over production height, and summing
incoherent produced-flavour contributions are all pipeline-level composition
concerns tied to this surface-to-detector step.

\begin{theorybox}
\subsection*{The production-flux-weighted height average, computed directly}
\pyfunc{propagate\_atmosphere\_grid\_to\_detector} (below) needs the same
production-flux-weighted probability average over production height that
\cref{ch:atmosphere} already derives for a single flavour
(\pyfunc{atmosphere\_probability\_integrated\_height}), but here for the
\emph{full} transition matrix $P_{\beta\alpha}(\theta,E,h)$ across every
zenith angle and produced flavour $\alpha$ at once. Rather than calling
that single-flavour function once per $(\theta,\alpha)$ pair, it evaluates
the same weighted-average identity directly on the stacked tensors,
\begin{equation}
  \keyresult{\langle P\rangle_h(\theta,E)_{\beta\alpha}
  = \frac{\int P_{\beta\alpha}(\theta,E,h)\,\phi^{\rm prod}_\alpha(\theta,E,h)\,\mathrm{d}h}
         {\int \phi^{\rm prod}_\alpha(\theta,E,h)\,\mathrm{d}h}},
\end{equation}
with the produced-flavour axis $\alpha$ broadcast against the transition
matrix's own two flavour indices via an extra singleton dimension on the
weight (\code{production\_flux.unsqueeze(-2)}), so a single pair of
trapezoidal integrals (numerator and denominator) produces the
height-averaged transition matrix for every $(\theta,\alpha,\beta)$
combination in one call. The corresponding height-\emph{integrated} (not
averaged) detector flux, $\Phi_\beta(\theta,E)=\sum_\alpha\int
P_{\beta\alpha}\,\phi^{\rm prod}_\alpha\,\mathrm{d}h$, is obtained from the
transition matrix directly via
\pyfunc{core.common.flux.flux\_transition}/\pyfunc{flux\_integrated\_coordinate}
(\cref{ch:core-common}), mirroring the single-flavour identity
$\langle P\rangle_h\,\Phi_h=\int P(h)\,\phi^{\rm prod}(h)\,\mathrm{d}h$
already established in \cref{ch:atmosphere}.
\end{theorybox}

\section{\texttt{pipeline.atmosphere\_earth} --- surface-to-detector and composition helpers}

\modulehead{pipeline/atmosphere\_earth.py}{\pyfunc{propagate\_surface\_to\_detector}
(one produced flavour, one zenith angle) and
\pyfunc{propagate\_atmosphere\_grid\_to\_detector} (every produced flavour,
every zenith angle, full transition matrix), plus the detector-composition
helpers \pyfunc{detector\_flux\_from\_production},
\pyfunc{integrate\_detector\_flux\_over\_height}, and
\pyfunc{sum\_detected\_flavours}.}

\begin{implbox}
\begin{itemize}
  \item \pyclass{AtmosphereEarthDetectorResult} (frozen dataclass): the
    input \code{surface} (\pyclass{AtmosphereSurfaceResult},
    \cref{ch:pipeline-atmosphere}), the Earth evolutor \code{S\_earth}, the
    composed \code{S\_total} ($S_{\rm earth}\,S_{\rm atmosphere}$),
    \code{detector\_states}/\code{detector\_probabilities} for the single
    produced flavour carried by \code{surface}, the full
    \code{transition\_probabilities} matrix, and -- when
    \code{surface.production} carries a height-resolved flux -- the same
    height-averaged/-integrated fields as \pyclass{AtmosphereSurfaceResult},
    now evaluated at the detector.
  \item \pyfunc{propagate\_surface\_to\_detector(surface, config, *,
    profile\_earth=None, integrate\_energy=False)}: builds (or reuses)
    the Earth profile, evaluates
    \pyfunc{earth\_evolutor\_from\_zenith} (\cref{ch:earth}) at the
    production zenith angle, composes it with \code{surface.S\_atmosphere},
    and applies the result to \code{surface.initial\_state}.
  \item \pyclass{AtmosphereDetectorGridResult} (frozen dataclass): the
    stacked, multi-angle, full-transition observables returned by the
    function below -- \code{transition\_probability\_theta\_Eh\_beta\_alpha},
    \code{probability\_height\_theta\_E\_beta\_alpha},
    \code{production\_flux\_theta\_Eh\_alpha},
    \code{detector\_flux\_theta\_Eh\_beta}, the height-, and optionally
    angle- and energy-integrated fluxes, plus the individual
    \code{detector\_results} (one \pyclass{AtmosphereEarthDetectorResult}
    per zenith angle).
  \item \pyfunc{propagate\_atmosphere\_grid\_to\_detector(production\_by\_flavour,
    config, *, flavour\_order=None, profile\_earth=None,
    trajectory\_steps=200, integrate\_angular=False, integrate\_energy=False)}:
    given \code{\{flavour: [production dict per zenith angle]\}} (aligned
    energy/height/angle grids across flavours, missing flavours treated as
    zero production flux), propagates each produced flavour at each zenith
    angle in turn, stacks the resulting transition matrices and production
    fluxes, and evaluates the weighted-average identity derived above.
  \item \pyfunc{detector\_flux\_from\_production(production\_flux,
    detector\_probabilities)} / \pyfunc{integrate\_detector\_flux\_over\_height(%
    h\_grid\_km, detector\_flux)} / \pyfunc{sum\_detected\_flavours(%
    integrated\_flux\_by\_production\_flavour)}: the three composition
    primitives \pyfunc{propagate\_atmosphere\_grid\_to\_detector} is built
    from, also usable directly by a caller assembling a custom multi-source
    composition -- respectively \pyfunc{core.common.flux.flux\_state}, a
    plain trapezoidal height integral, and an incoherent sum over a
    \code{\{flavour: flux\}} mapping (requiring at least one entry and
    matching shapes across all of them).
\end{itemize}
\end{implbox}

\begin{examplebox}[title={Continuing the atmosphere surface state through the Earth}]
\begin{lstlisting}[style=tpython]
from tpeanuts.pipeline.atmosphere_earth import propagate_surface_to_detector

# `surface` from the pipeline.atmosphere example (Chapter 13): a coherent
# nu_mu state produced at theta_deg=60, already propagated to the surface.
detector = propagate_surface_to_detector(surface, config)

print(detector.detector_probabilities.shape)     # (E, h, flavour)
print(detector.transition_probabilities.shape)   # (E, h, flavour_out, flavour_in)
\end{lstlisting}
\end{examplebox}

Summing several produced-flavour contributions at the detector, once each
has already been propagated and integrated over production height, is a
plain incoherent sum over a name $\to$ tensor mapping:
\begin{examplebox}[title={Combining $\nu_\mu$- and $\nu_e$-produced contributions at the detector}]
\begin{lstlisting}[style=tpython]
from tpeanuts.pipeline.atmosphere_earth import (
    integrate_detector_flux_over_height,
    sum_detected_flavours,
)

h_km = surface.production["h_grid_km"]
numu_integrated = integrate_detector_flux_over_height(h_km, detector.detector_flux_Eh)
nue_integrated = integrate_detector_flux_over_height(h_km, nue_detector.detector_flux_Eh)

total_detector_flux = sum_detected_flavours({
    "numu": numu_integrated,
    "nue": nue_integrated,
})
\end{lstlisting}
\end{examplebox}
